\documentclass[11pt]{article}

\usepackage[T1]{fontenc}
\usepackage[utf8]{inputenc}
\usepackage{amsmath}
\usepackage{amssymb}
\usepackage{geometry}
\usepackage{listings}
\usepackage{xcolor}

\geometry{margin=1in}

\definecolor{codebg}{RGB}{245,245,245}
\definecolor{codecomment}{RGB}{80,80,80}

\lstset{
  basicstyle=\ttfamily\small,
  backgroundcolor=\color{codebg},
  commentstyle=\color{codecomment},
  frame=single,
  columns=fullflexible,
  keepspaces=true,
  showstringspaces=false
}

\title{Deriving the CRC-8 Update From $R_n$ to $R_{n+1}$}
\author{}
\date{}

\begin{document}
\maketitle

\section{Goal}

For a serial, MSB-first CRC, the update logic is often written in hardware as
\[
\mathit{fb} = \mathit{mosi} \oplus \mathit{crc}[7].
\]
This can look surprising at first, because long division is usually described as
``look at the current MSB, then divide.'' The missing detail is that the next
dividend bit is brought down at the same step. In the LFSR form, those two actions
collapse into the single feedback term above.

This note derives the recurrence for the next remainder $R_{n+1}$ from the current
remainder $R_n$ and then maps the result directly to Verilog.

\section{CRC Definition}

Let the generator polynomial be
\[
g(x) = x^8 + p(x),
\]
where $p(x)$ contains the lower-order tap terms.

For CRC-8/AUTOSAR,
\[
g(x) = x^8 + x^4 + x^3 + x^2 + 1,
\]
so
\[
p(x) = x^4 + x^3 + x^2 + 1.
\]

Let $D_n(x)$ be the polynomial formed by the first $n$ message bits that have
already been processed, and let $R_n(x)$ be the current 8-bit remainder. The
standard CRC definition is
\[
R_n(x) = x^8 D_n(x) \bmod g(x).
\]

\section{Appending One More Bit}

Suppose the next input bit is $b_n \in \{0,1\}$. Appending one bit to the message
polynomial gives
\[
D_{n+1}(x) = xD_n(x) + b_n.
\]

Substitute that into the CRC definition:
\begin{align*}
R_{n+1}(x)
&= x^8 D_{n+1}(x) \bmod g(x) \\
&= x^8 (xD_n(x) + b_n) \bmod g(x) \\
&= (x \cdot x^8 D_n(x) + b_n x^8) \bmod g(x).
\end{align*}

Because $x^8 D_n(x) \bmod g(x) = R_n(x)$, this becomes
\[
R_{n+1}(x) = \bigl(xR_n(x) + b_n x^8\bigr) \bmod g(x).
\]

\section{Where the Feedback XOR Comes From}

Write the current remainder as
\[
R_n(x) = r_7 x^7 + r_6 x^6 + r_5 x^5 + r_4 x^4 + r_3 x^3 + r_2 x^2 + r_1 x + r_0.
\]

After multiplying by $x$,
\[
xR_n(x) = r_7 x^8 + r_6 x^7 + r_5 x^6 + r_4 x^5 + r_3 x^4 + r_2 x^3 + r_1 x^2 + r_0 x.
\]

Now include the new message bit:
\[
xR_n(x) + b_n x^8 =
(r_7 \oplus b_n)x^8 + r_6 x^7 + r_5 x^6 + r_4 x^5 + r_3 x^4 + r_2 x^3 + r_1 x^2 + r_0 x.
\]

The coefficient of $x^8$ is therefore
\[
f = r_7 \oplus b_n.
\]

This is exactly the serial hardware feedback term:
\[
\boxed{f = \mathit{crc}[7] \oplus \mathit{mosi}}
\]
for an MSB-first implementation.

\section{Reducing the $x^8$ Term}

Since
\[
g(x) = x^8 + p(x),
\]
we have the congruence
\[
x^8 \equiv p(x) \pmod{g(x)}.
\]

So every time the coefficient $f$ of $x^8$ is 1, the polynomial $p(x)$ must be XORed
back into the shifted remainder. That yields the recurrence
\[
R_{n+1}(x) = \operatorname{shift}(R_n(x)) \oplus f\,p(x),
\]
where
\[
\operatorname{shift}(R_n(x)) = r_6 x^7 + r_5 x^6 + r_4 x^5 + r_3 x^4 + r_2 x^3 + r_1 x^2 + r_0 x.
\]

So the generic serial CRC update is
\[
\boxed{R_{n+1} = \operatorname{shift}(R_n) \oplus (r_7 \oplus b_n)\,p(x)}.
\]

\section{Bit Equations for CRC-8/AUTOSAR}

For CRC-8/AUTOSAR,
\[
p(x) = x^4 + x^3 + x^2 + 1.
\]

That means the feedback bit $f$ is XORed into the stages corresponding to
$x^4$, $x^3$, $x^2$, and $x^0$. The next-state equations are
\begin{align*}
r_7' &= r_6 \\
r_6' &= r_5 \\
r_5' &= r_4 \\
r_4' &= r_3 \oplus f \\
r_3' &= r_2 \oplus f \\
r_2' &= r_1 \oplus f \\
r_1' &= r_0 \\
r_0' &= f.
\end{align*}

This is the usual MSB-first, serial LFSR form.

\section{Matching Verilog}

\begin{lstlisting}[language=Verilog]
module crc8_autosar_serial (
  input  wire       clk,
  input  wire       rst_n,
  input  wire       enable,
  input  wire       mosi,
  output reg  [7:0] crc
);

  wire fb = crc[7] ^ mosi;

  always @(posedge clk or negedge rst_n) begin
    if (!rst_n) begin
      crc <= 8'hFF; // AUTOSAR seed
    end else if (enable) begin
      crc[7] <= crc[6];
      crc[6] <= crc[5];
      crc[5] <= crc[4];
      crc[4] <= crc[3] ^ fb;
      crc[3] <= crc[2] ^ fb;
      crc[2] <= crc[1] ^ fb;
      crc[1] <= crc[0];
      crc[0] <= fb;
    end
  end
endmodule
\end{lstlisting}

If the final AUTOSAR CRC byte is needed, apply the final XOR after all bits are
processed:
\[
\mathit{crc\_out} = \mathit{crc} \oplus \texttt{8'hFF}.
\]

\section{Why \texttt{mosi} Must Be in the Feedback}

If the update used only $\mathit{crc}[7]$, then the state evolution would depend
only on the previous remainder, not on the actual incoming message bit at every
tap point required by the polynomial. That would no longer implement the same
polynomial division.

The feedback XOR
\[
\mathit{crc}[7] \oplus \mathit{mosi}
\]
is the coefficient of the new $x^8$ term after both of the following happen:
\begin{enumerate}
\item the previous remainder is shifted by one power of $x$;
\item the next message bit is appended.
\end{enumerate}

That is why the incoming bit affects the same reduction decision as the outgoing
MSB of the remainder.

\section{Alternate Derivation Using Physical Bit Insertion}

The previous derivation used the standard CRC definition
\[
R_n(x) = x^8 D_n(x) \bmod g(x).
\]
It is also possible to derive a valid serial update rule directly from the physical
LFSR behavior, without starting from that definition.

Assume the register is interpreted as
\[
S_n(x) = s_0 + s_1 x + s_2 x^2 + \cdots + s_7 x^7,
\]
where $s_i = \texttt{crc}[i]$.

Now assume the serial hardware behaves as follows:
\begin{enumerate}
\item the new input bit $m$ (that is, \texttt{mosi}) is inserted into stage \texttt{crc[0]};
\item all existing register contents shift left by one stage on the clock edge.
\end{enumerate}

Under that convention, the raw shift-and-insert step is
\[
T(x) = xS_n(x) + m.
\]

This is because multiplying by $x$ shifts every stored term up by one power, while
adding $m$ inserts the new bit at $x^0$.

Expanding,
\[
T(x) = m + s_0 x + s_1 x^2 + s_2 x^3 + s_3 x^4 + s_4 x^5 + s_5 x^6 + s_6 x^7 + s_7 x^8.
\]

The only term that does not fit in an 8-bit register is the overflow term $s_7 x^8$.
For CRC-8/AUTOSAR,
\[
g(x) = x^8 + x^4 + x^3 + x^2 + 1.
\]
Working modulo $g(x)$ means that $g(x)$ itself is equivalent to zero:
\[
g(x) \equiv 0 \pmod{g(x)}.
\]
Therefore,
\[
x^8 + x^4 + x^3 + x^2 + 1 \equiv 0 \pmod{g(x)}.
\]
In $\mathrm{GF}(2)$, subtraction and addition are the same operation, so rearranging
the terms gives
\[
x^8 \equiv x^4 + x^3 + x^2 + 1 \pmod{g(x)}.
\]
This is why the overflow term is reduced as
\[
s_7 x^8 \equiv s_7(x^4 + x^3 + x^2 + 1).
\]

Substituting that back in,
\[
S_{n+1}(x) =
m + s_0 x + s_1 x^2 + s_2 x^3 + s_3 x^4 + s_4 x^5 + s_5 x^6 + s_6 x^7
+ s_7(x^4 + x^3 + x^2 + 1).
\]

Collecting powers of $x$ gives
\[
S_{n+1}(x) =
(m \oplus s_7)
+ s_0 x
+ (s_1 \oplus s_7)x^2
+ (s_2 \oplus s_7)x^3
+ (s_3 \oplus s_7)x^4
+ s_4 x^5
+ s_5 x^6
+ s_6 x^7.
\]

Therefore the next-state equations are
\begin{align*}
s_0' &= m \oplus s_7 \\
s_1' &= s_0 \\
s_2' &= s_1 \oplus s_7 \\
s_3' &= s_2 \oplus s_7 \\
s_4' &= s_3 \oplus s_7 \\
s_5' &= s_4 \\
s_6' &= s_5 \\
s_7' &= s_6.
\end{align*}

So \texttt{mosi} is still part of the update in this convention, but it enters only
through the newly inserted low-order term:
\[
s_0' = \texttt{mosi} \oplus s_7.
\]
The tapped reductions at $x^2$, $x^3$, and $x^4$ come only from reducing the
overflow term $s_7 x^8$, so they depend on $s_7$ alone:
\[
s_2' = s_1 \oplus s_7,\qquad
s_3' = s_2 \oplus s_7,\qquad
s_4' = s_3 \oplus s_7.
\]
That is why this alternate convention does not collapse into the same single
shared feedback term $\texttt{crc[7]} \oplus \texttt{mosi}$ used by the earlier
MSB-first AUTOSAR derivation.

The corresponding Verilog is
\begin{lstlisting}[language=Verilog]
module crc8_shift_left_insert_lsb (
  input  wire       clk,
  input  wire       rst_n,
  input  wire       enable,
  input  wire       mosi,
  output reg  [7:0] crc
);

  wire overflow = crc[7];

  always @(posedge clk or negedge rst_n) begin
    if (!rst_n) begin
      crc <= 8'h00;
    end else if (enable) begin
      crc[0] <= mosi ^ overflow;
      crc[1] <= crc[0];
      crc[2] <= crc[1] ^ overflow;
      crc[3] <= crc[2] ^ overflow;
      crc[4] <= crc[3] ^ overflow;
      crc[5] <= crc[4];
      crc[6] <= crc[5];
      crc[7] <= crc[6];
    end
  end
endmodule
\end{lstlisting}

\section{Relation Between the Two Derivations}

This alternate derivation is mathematically consistent, but it is a different
serial convention from the earlier MSB-first AUTOSAR form.

In the earlier derivation,
\begin{itemize}
\item the state is interpreted as the standard CRC remainder;
\item the feedback term is $\texttt{crc[7]} \oplus \texttt{mosi}$;
\item the result matches the usual non-reflected AUTOSAR CRC-8 implementation.
\end{itemize}

In the alternate derivation above,
\begin{itemize}
\item the state is interpreted directly from the physical register layout;
\item the new bit is inserted at $x^0$ and then shifted toward higher powers;
\item \texttt{mosi} affects the low-order update directly, while the tap XORs come
from reducing the overflow term $s_7 x^8$;
\item the resulting recurrence is different, even though it still comes from the
same generator polynomial.
\end{itemize}

So the placement of the incoming bit at the physical low end of the register does
not by itself determine the CRC convention. The state definition and shift
direction determine which recurrence is correct.

\end{document}
